% answers_short.tex —— 1988 年·广东卷·文史类 · 简短答案速查
% 由 tools/make_short_answers.py 自 answers.tex 的 \gksolhead/\gkanswer 提取，
% 并与 1988/1988_数学真题_广东文_答案要点.md 逐题比对后生成，勿手改（改源 MD 后重跑）。

\begin{gknotice}
\textbf{使用说明}：本节只给最终答案便于快速核对，完整推导见解析版。答案为编写者独立推导，\textbf{非官方标准答案}。
\end{gknotice}

\gkpart{大题}
\begin{itemize}[leftmargin=2.6em]
  \item[\textbf{1.}] (1) B；(2) B；(3) B；(4) B；(5) B；(6) C；(7) D；(8) D；(9) C；(10) A；(11) B；(12) C；(13) B；(14) B；(15) C；(16) B；(17) D；(18) D；(19) D；(20) B；(21) C；(22) C；(23) A；(24) B；(25) C
  \item[\textbf{2.}] (1) $y=-x+2$；(2) $\dfrac{\sqrt2}{2}$；(3) $\pi$；(4) $\sqrt3-i$；(5) $-1$
  \item[\textbf{3.}] (1) $|u|=2\sqrt2$，$\arg u=\dfrac{3\pi}{2}$；(2) $x=8,\ y=-8$
  \item[\textbf{4.}] (1) $DA\perp AB$；(2) $16\,\mathrm{cm}^3$
  \item[\textbf{5.}] (1) $\{a_n\}$ 不是等差数列；(2) 常数为 $\dfrac32$
  \item[\textbf{6.}] (1) $\dfrac1{r^2}-k^2=1$，且 $|k|\ne1$；(2) 不能经过双曲线的任一焦点
\end{itemize}
